\chapter{Clase 16: Cuadraturas del problema relativo de dos cuerpos}

Esta clase continúa el desarrollo analítico del problema relativo de dos cuerpos introducido en la clase anterior. El objetivo es identificar las \textbf{integrales primeras} o \textbf{cuadraturas} que permiten resolver completamente la dinámica relativa y establecer por qué este problema es integrable de forma exacta.

\section{Planteamiento del sistema y necesidad de cuadraturas}
Partimos de la ecuación diferencial que describe el movimiento relativo entre dos cuerpos sometidos a gravitación mutua:
\[
\ddot{\vec{r}} = -\frac{\mu}{r^3}\vec{r}
\]
donde \(\mu = G(m_1 + m_2)\) y \(r = |\vec{r}|\).

Esta es una ecuación diferencial vectorial de segundo orden. Al descomponerla en coordenadas cartesianas, el problema queda gobernado por \textbf{6 variables dinámicas independientes} \((x, y, z, \dot{x}, \dot{y}, \dot{z})\). Para integrar completamente el sistema y obtener una solución cerrada, es necesario hallar \textbf{6 cuadraturas} o integrales primeras independientes.

Aunque el problema de \(N\) cuerpos posee constantes de movimiento globales, el problema relativo de dos cuerpos constituye un sistema autónomo nuevo que requiere su propia deducción analítica paso a paso.

\begin{importante}
La meta de esta clase no es todavía escribir la órbita final en forma cónica, sino mostrar cómo la dinámica admite suficientes constantes de movimiento para cerrar completamente el sistema.
\end{importante}

\section{Primera cuadratura: momento angular relativo específico}
\subsection{Obtención mediante factor integrante}
Multiplicamos vectorialmente la ecuación de movimiento por el vector posición \(\vec{r}\) por la izquierda:
\[
\vec{r} \times \ddot{\vec{r}} = -\frac{\mu}{r^3}(\vec{r} \times \vec{r}) = \vec{0}
\]

Aplicamos la regla de derivación del producto vectorial:
\[
\frac{d}{dt}(\vec{r} \times \dot{\vec{r}}) = \dot{\vec{r}} \times \dot{\vec{r}} + \vec{r} \times \ddot{\vec{r}} = \vec{0}
\]

Al integrar respecto al tiempo, obtenemos una constante vectorial:
\[
\vec{h} \equiv \vec{r} \times \vec{v} = \text{cte.}
\]

Esta integral aporta \textbf{3 cuadraturas} (una por cada componente cartesiana). El vector \(\vec{h}\) es el \textbf{momento angular relativo específico}, es decir, por unidad de masa reducida.

\subsection{Interpretación física y geométrica}
\begin{itemize}
    \item \textbf{Movimiento plano:} como \(\vec{h}\) es constante y perpendicular tanto a \(\vec{r}\) como a \(\vec{v}\), el movimiento relativo queda confinado a un plano fijo perpendicular a \(\vec{h}\).
    \item \textbf{Coordenadas polares:} en el plano orbital, su módulo se expresa como \(h = r^2\dot{\theta}\).
    \item \textbf{Velocidad areal y segunda ley de Kepler:} el área infinitesimal barrida por el radio vector es \(dA = \frac{1}{2}r^2 d\theta\). Derivando respecto al tiempo:
\end{itemize}

\[
\frac{dA}{dt} = \frac{1}{2}r^2\dot{\theta} = \frac{h}{2} = \text{cte.}
\]

Esto demuestra que la velocidad areal es constante. La magnitud de \(h\) controla directamente la tasa de barrido del área, lo cual constituye la formulación dinámica de la segunda ley de Kepler.

\begin{importante}
La conservación del momento angular específico no solo reduce el número de grados efectivos del problema, sino que además garantiza que toda órbita kepleriana es plana.
\end{importante}

\section{Segunda cuadratura: energía relativa específica}
\subsection{Obtención mediante producto escalar}
Multiplicamos escalarmente la ecuación de movimiento por la velocidad \(\dot{\vec{r}}\):
\[
\dot{\vec{r}} \cdot \ddot{\vec{r}} = -\frac{\mu}{r^3}\vec{r} \cdot \dot{\vec{r}}
\]

Identificamos las derivadas exactas:
\[
\frac{1}{2}\frac{d}{dt}(v^2) = -\mu \frac{d}{dt}\left(\frac{1}{r}\right)
\]

Reordenando e integrando:
\[
\frac{d}{dt}\left(\frac{v^2}{2} - \frac{\mu}{r}\right) = 0
\]

Definimos la \textbf{energía relativa específica} como la constante resultante:
\[
\mathcal{E} \equiv \frac{v^2}{2} - \frac{\mu}{r} = \text{cte.}
\]

Esta integral escalar aporta \textbf{1 cuadratura} adicional.

\subsection{Interpretación física}
Representa la conservación de la energía mecánica por unidad de masa reducida. Su valor algebraico clasifica la trayectoria:
\begin{itemize}
    \item \(\mathcal{E} < 0\): órbita ligada (elipse).
    \item \(\mathcal{E} = 0\): caso límite de escape (parábola).
    \item \(\mathcal{E} > 0\): órbita no ligada (hipérbola).
\end{itemize}

\begin{importante}
El signo de \(\mathcal{E}\) permite distinguir si el movimiento permanece acotado o si el cuerpo escapa al infinito, incluso antes de resolver explícitamente la trayectoria.
\end{importante}

\section{Tercera cuadratura: vector de excentricidad}
\subsection{Construcción del factor integrante}
Para completar el sistema, cruzamos la ecuación de movimiento con el vector constante \(\vec{h}\):
\[
\ddot{\vec{r}} \times \vec{h} = -\frac{\mu}{r^3}\vec{r} \times (\vec{r} \times \vec{v})
\]

Aplicamos la identidad del triple producto vectorial (BAC-CAB):
\[
\vec{r} \times (\vec{r} \times \vec{v}) = \vec{r}(\vec{r} \cdot \vec{v}) - \vec{v}(\vec{r} \cdot \vec{r}) = r\dot{r}\vec{r} - r^2\vec{v}
\]

Sustituyendo y simplificando:
\[
\ddot{\vec{r}} \times \vec{h} = \mu\left(\frac{\vec{v}}{r} - \frac{\dot{r}\vec{r}}{r^2}\right)
\]

El término entre paréntesis coincide con la derivada temporal del vector unitario radial \(\hat{r} = \vec{r}/r\):
\[
\frac{d}{dt}\left(\frac{\vec{r}}{r}\right) = \frac{\dot{\vec{r}}}{r} - \frac{\vec{r}\dot{r}}{r^2}
\]

Por tanto, la ecuación se reduce a una derivada total exacta:
\[
\frac{d}{dt}(\dot{\vec{r}} \times \vec{h}) = \frac{d}{dt}\left(\mu\frac{\vec{r}}{r}\right)
\]

Integrando, obtenemos una nueva constante vectorial:
\[
\vec{e} \equiv \frac{\vec{v} \times \vec{h}}{\mu} - \frac{\vec{r}}{r} = \text{cte.}
\]

\subsection{Propiedades del vector de excentricidad}
Este vector, también llamado \textbf{vector de Laplace} o \textbf{vector de Runge--Lenz}, posee varias propiedades fundamentales:
\begin{itemize}
    \item yace en el plano orbital,
    \item es perpendicular a \(\vec{h}\),
    \item apunta hacia el periastro,
    \item su módulo coincide con la excentricidad orbital: \(e = |\vec{e}|\).
\end{itemize}

Formalmente aporta \textbf{3 cuadraturas}, aunque no todas son independientes de las integrales anteriores.

\begin{importante}
El vector de excentricidad contiene información geométrica fina sobre la orientación y la forma de la órbita. Mientras \(\vec{h}\) fija el plano orbital, \(\vec{e}\) fija la dirección del periastro dentro de ese plano.
\end{importante}

\section{Cierre del sistema: conteo de cuadraturas independientes}
Sumando las constantes obtenidas, tenemos:
\begin{itemize}
    \item momento angular relativo específico \(\vec{h}\): 3 componentes,
    \item energía relativa específica \(\mathcal{E}\): 1 escalar,
    \item vector de excentricidad \(\vec{e}\): 3 componentes.
\end{itemize}

El total aparente es de \(7\) ecuaciones. Sin embargo, el sistema no posee siete grados de libertad independientes, porque existen relaciones algebraicas que ligan estas constantes:
\begin{enumerate}
    \item ortogonalidad geométrica:
\end{enumerate}
\[
\vec{h} \cdot \vec{e} = 0
\]
\begin{enumerate}[resume]
    \item relación escalar derivada de \(\vec{e} \cdot \vec{e}\):
\end{enumerate}
\[
 e^2 = 1 + \frac{2\mathcal{E}h^2}{\mu^2}
\]

Al restar estas dependencias funcionales, el sistema queda con exactamente \textbf{6 cuadraturas independientes}, que coinciden con el número de variables dinámicas iniciales. Esto demuestra que el problema relativo de dos cuerpos es \textbf{completamente integrable} y permite derivar la ecuación de la trayectoria \(r(\theta)\) sin recurrir, en principio, a métodos numéricos.

\section{Síntesis}
En esta clase se obtuvo el conjunto fundamental de integrales primeras del problema relativo de dos cuerpos:
\begin{itemize}
    \item conservación del momento angular específico,
    \item conservación de la energía específica,
    \item conservación del vector de excentricidad.
\end{itemize}

Estas cuadraturas no solo permiten cerrar analíticamente el sistema, sino que preparan el camino para la deducción explícita de la ecuación de la cónica orbital y el estudio detallado de los elementos orbitales en la clase siguiente.